\documentclass[../main.tex]{subfiles}
\ifSubfilesClassLoaded{\setcounter{chapter}{10}}{}
\begin{document}

\chapter{Comprehension, Generator, dan Iterator Pengenalan}

\begin{subcpmk}
  \item Sub-CPMK 7.2: Menulis list/dict/set comprehension dan fungsi generator sederhana dengan \texttt{yield}\newline untuk menghasilkan deret data secara hemat memori
\end{subcpmk}

\SesiRPS{12}{7.2}{$\sim$170 menit}{CPMK-7}

\section{Sarana dan prasarana}

Python 3.10+, editor, terminal; tidak memerlukan paket pihak ketiga.

\section{Keselamatan kerja dan etika}

Etika akademik: comprehension dan generator harus Anda pahami, bukan hanya disalin tanpa memahami alur data.

\section{Prosedur kerja}

\begin{enumerate}[leftmargin=*]
  \item Pelajari sintaks comprehension dari contoh.
  \item Bandingkan dengan loop klasik untuk kasus yang sama.
  \item Tulis satu generator dengan \texttt{yield}; iterasikan dengan \texttt{for} atau \texttt{next()}.
  \item Refaktor potongan kode loop panjang menjadi gaya idiomatik (aktivitas).
\end{enumerate}

\section{List dan Dict Comprehension}

\begin{lstlisting}[caption=Comprehension]
squares = [x * x for x in range(10) if x % 2 == 0]
nama_ke_id = {n.lower(): i for i, n in enumerate(["Ali", "Budi", "Ali"])}
\end{lstlisting}

\section{Generator dengan yield}

\begin{lstlisting}[caption=Generator sederhana]
def countdown(n):
    while n >= 0:
        yield n
        n -= 1

for x in countdown(3):
    print(x)
\end{lstlisting}

\begin{aktivitas}
  \item Dari list bilangan \texttt{data}, buat list baru berisi hanya bilangan positif ganjil memakai list comprehension.
  \item Buat dict comprehension yang memetakan huruf pertama kata ke panjang kata dari daftar string.
  \item Tulis generator \texttt{bilangan\_genap(maks)} yang menghasilkan bilangan genap dari 0 sampai \texttt{maks} (inklusif), lalu cetak hasilnya.
  \item Simpan pekerjaan di \texttt{NIM\_TIF105\_M12\_comp.py}.
\end{aktivitas}

\begin{latihan}
  \item Kapan generator lebih tepat daripada list comprehension?
  \item Apa perilaku iterator setelah habis?
\end{latihan}

\begin{asesmen}
\textbf{Formatif:} comprehension valid secara sintaks; generator dapat diiterasi; refaktor mengurangi loop bersarang berlebihan tanpa mengubah hasil logika.
\end{asesmen}

\begin{checklist}
  \item Saya dapat menulis list/dict comprehension dasar
  \item Saya dapat membuat fungsi generator dengan \texttt{yield}
  \item Saya memahami kapan memakai generator untuk efisiensi memori
\end{checklist}

\begin{rangkuman}
Comprehension memadatkan loop pembangun koleksi; generator menunda komputasi dengan \texttt{yield}.

\textbf{Kata kunci:} comprehension, generator, \code{yield}, iterator
\end{rangkuman}

\section{Referensi sesi}

\begin{itemize}[leftmargin=*]
  \item \textbf{[14]} Pilgrim, M. \textit{Dive Into Python 3} --- iterator dan generator.
  \item \textbf{[8, 9]} \textit{The Python Tutorial} --- generator expressions.
  \item \textbf{[19]} Real Python --- comprehension dan generator.
\end{itemize}

\end{document}
